\documentclass[12pt]{article}
\usepackage[utf8]{inputenc}
\usepackage[T1]{fontenc}
\usepackage[spanish]{babel}
\begin{document}
\title{Proyecto 5 — Opción Americana · Árbol Binomial}
\maketitle

Proyecto 5 — Opción Americana · Árbol Binomial
Finanzas Cuantitativas · Proyecto 5
Opción Americana
mediante Árbol Binomial CRR
¿Cómo cambia la valuación cuando el ejercicio puede ocurrir
antes
del vencimiento? Comparativa PUT Europea vs. Americana
Marco Teórico
Modelo CRR — Cox, Ross & Rubinstein (1979)
El modelo binomial divide el tiempo al vencimiento
T
en
N
pasos iguales de longitud
Δt = T/N
. En cada paso el activo puede subir un factor
u
o bajar un factor
d
.
La clave está en la
inducción hacia atrás (backward induction)
: se calculan los payoffs en el vencimiento y se descuenta hacia el presente, pero en cada nodo intermedio se verifica si conviene ejercer anticipadamente (sólo para opciones americanas).
Parámetros del árbol
\[
          u = e^{\sigma\sqrt{\Delta t}}, \quad d = e^{-\sigma\sqrt{\Delta t}} = \frac{1}{u}
        \]
\[
          p = \frac{e^{r\Delta t} - d}{u - d}
        \]
donde
p
es la
probabilidad neutral al riesgo
de movimiento alcista.
Backward Induction
\[
          \text{Payoff: } \max(K - S, 0)
        \]
\[
          V_{i,j} = e^{-r\Delta t}\bigl[p\,V_{i+1,j+1} + (1-p)\,V_{i+1,j}\bigr]
        \]
Para la
opción americana
se agrega:
\[
          V_{i,j}^{\text{Am}} = \max\!\Bigl(K - S_{i,j},\; V_{i,j}^{\text{cont}}\Bigr)
        \]
Diferencia fundamental: PUT Europea vs. Americana
La
PUT Europea
sólo puede ejercerse al vencimiento. La
PUT Americana
puede ejercerse en cualquier momento, lo que añade valor cuando el activo cae profundamente por debajo del strike, ya que el inversor prefiere recibir el beneficio
hoy
en lugar de esperar.
Por ello:
\(P_{\text{Am}} \geq P_{\text{Eu}}\)
siempre. La diferencia se llama
Prima por Ejercicio Temprano (PET)
.
Parámetros del Modelo
⬧ Modifica los parámetros y presiona Recalcular
S₀ — Precio inicial
K — Strike
r — Tasa libre de riesgo
σ — Volatilidad
T — Tiempo (años)
N — Pasos del árbol
Recalcular
Resultados
Put Europea
—
Valor sólo al vencimiento
Put Americana
—
Ejercicio anticipado incluido
Prima por Ej. Temprano
—
Americana − Europea
Nodos con ej. anticipado
—
Nodos en rojo en el árbol
Factor u (alza)
—
Factor d (baja)
—
Prob. neutral p
—
Δt (años/paso)
—
Árbol Binomial — Precios del Activo
Precio del activo S
i,j
Nodo intermedio
Nodo terminal (vencimiento)
Árbol Binomial — Valor de la Opción
▸ PUT EUROPEA
▸ PUT AMERICANA
Ejercicio anticipado óptimo
Continuar (no ejercer)
Terminal
Comparativa por Precios del Activo
La brecha entre ambas curvas es la
Prima por Ejercicio Temprano
Cálculos Paso a Paso
Presiona
Recalcular
para ver el proceso detallado…
Tabla Comparativa
Característica
Put Europea
Put Americana
Ejercicio
Sólo al vencimiento
En cualquier momento
Valuación (árbol)
—
—
Curva del árbol
Backward induction puro
Backward + max(payoff, cont.)
Intrinsic value posible
Sólo en \(t=T\)
En cualquier nodo
Precio siempre ≥
—
Sí, ≥ Europea
Prima por ej. temprano
0.0000
—
Nodos de ej. temprano
N/A
—
Cuándo se prefiere ejercer
Nunca antes
Cuando \(K - S > V_{\text{cont}}\)
Conclusiones
01
¿Por qué la opción americana vale más?
La opción americana otorga al tenedor
todos los derechos
de la europea más la flexibilidad de ejercer anticipadamente. Este derecho adicional tiene valor positivo: el inversor puede capturar ganancias inmediatas cuando el activo cae profundamente y recibir el beneficio
hoy
en lugar de esperar, aprovechando el valor del dinero en el tiempo.
02
¿Cuándo conviene ejercer temprano?
Conviene ejercer anticipadamente cuando el
valor intrínseco
\((K - S)\) supera el
valor de continuación
. Esto ocurre cuando el activo cae muy por debajo del strike, la tasa de interés es alta (el valor presente del strike es mayor hoy) y el tiempo restante no compensa la espera.
03
La Prima por Ejercicio Temprano (PET)
La PET = \(P_{\text{Am}} - P_{\text{Eu}} \geq 0\) representa el valor monetario de la flexibilidad. Crece con la volatilidad (más nodos donde puede ser óptimo ejercer), con la tasa libre de riesgo y cuando el activo está profundamente
in-the-money
. Es el precio que el mercado asigna a la opción adicional de ejercicio anticipado.
04
Implicaciones para la cobertura (hedging)
Un vendedor de una put americana debe considerar la PET al cotizar el precio, pues asume el riesgo de asignación prematura. El tenedor debería ejercer cuando el nodo aparece en rojo en el árbol, señal de que continuar la opción destruye valor. El árbol binomial es la herramienta estándar en la industria para valorar derivados con esta característica.
Resumen del modelo con los parámetros actuales
Ejecuta el modelo para ver el resumen personalizado…
Proyecto 5 · Finanzas Cuantitativas · Modelo CRR (Cox–Ross–Rubinstein, 1979) · HTML puro + Chart.js + MathJax

\end{document}
